\chapter{Reference-Informed Judge Generation}\label{app:judge-generation}

This appendix details the rubric-generation mechanics referenced in
Chapter~\ref{ch:experimental-design}, \S\ref{sec:exp-rosters}.

\section{Generation Procedure}

The judge rubric generation procedure (Section~\ref{sec:arch-judge}) is
instantiated as follows. Representative-query packs are drawn from
each leaf's construction pool at pack size $k = 25$; a leaf with fewer
than 25 construction queries uses its full pool as a single pack.
The pool a pack is drawn from is the leaf's construction queries
\emph{after} two filters: the reserved-set filter, which is audited by
a build-failing assertion and is what makes the answer-key-blind claim
enforceable rather than asserted, and a hash filter that thins the pool
by roughly a fifth. A leaf's rubric is therefore induced from a sample
of its construction queries rather than from all of them.
The map-step LLM and the reduce-step LLM are the same underlying model
as the arena judge, to avoid introducing a second
model's stylistic biases into the rubric; this choice is revisited as
a limitation in Section~\ref{sec:disc-debt}. Guideline generation is
a one-time cost per frozen snapshot: rubrics are computed once at
snapshot-freeze time and reused for every subsequent arena run against
that snapshot, regardless of which models are being compared.

\section{Structural Template}\label{sec:judge-template-full}

The global structural template, covering presentation order, tie
handling, and output format, is fixed across all leaves. It governs
only the shared scaffolding of the judge prompt; the leaf-specific
reference-informed content (representative queries, answer-key
context) is injected into this template at generation time and does
not modify its structure.

% Full verbatim template goes here — insert the JudgePrompts.kt /
% GenericPairwiseJudgePrompt.kt template text (structural scaffold,
% tie-handling instructions, output-format specification) as a
% \begin{quote} or \lstlisting block once finalized from source.

\section{Dual-Call Position-Reversal Protocol}

Each logical comparison is realised as two judge calls with response
position reversed between calls, per the dual-call protocol of
Section~\ref{sec:arch-judge}. Disagreement between the two calls is
resolved as a tie, with each candidate credited half a win
(\parencite{davidson1970ties}), rather than resolved by a coin flip or
discarded.
